% GENERATED FILE. Edit canonical lessons, then run npm run generate:books.
% canonical-source-hash: fnv1a64:91ea89557074098c
% canonical-lessons: SA-C109-dadimam, SA-C109-draksa, SA-C109-panasam, SA-C109-nimbukam, SA-C109-apupah, SA-R109-first-pass-tools-and-things-you-use, SA-R109-second-pass-machines-the-house-and-food

\chapter{Pomegranate, Grapes, Jackfruit, Lemon, Cake}
\label{ch:sa-pomegranate-grapes-jackfruit-lemon-cake}
\hlchaptermodality{109}

\begin{chapteropening}
\textbf{By the end of this chapter:} \emph{I can say a pomegranate, grapes, a jackfruit, a lemon and a cake.}

Complete the last lesson of chapter 109: I can say a pomegranate, grapes, a jackfruit, a lemon and a cake.
\end{chapteropening}

\section[\texorpdfstring{\sk{दाडिमम्} (dāḍimam)}{दाडिमम् (dāḍimam)}]{\sk{दाडिमम्} (dāḍimam) — a pomegranate}

\label{lesson:SA-C109-dadimam}



\begin{quote}
Before the new one: say the Sanskrit for a pumpkin, then the Sanskrit for a coconut.
\end{quote}

\subsection*{You'll want to know: \sk{दाडिमम्}}
\textbf{\sk{दाडिमम्}} — \emph{dāḍimam} — \textquotedblleft{}a pomegranate\textquotedblright{}.

\begin{grammarlens}[title={What you've built}]
One more word for this chapter.
\end{grammarlens}

\subsection*{Your turn}
\begin{itemize}
\raggedright
  \item \emph{Say it:} \emph{dāḍimam}
  \item \emph{Say it:} \emph{dāḍimam}, once more
  \item \emph{Say it:} say \emph{dāḍimam}
  \item \emph{Recall:} say the Sanskrit for a pumpkin, then the Sanskrit for a coconut, then say \emph{dāḍimam} again
\end{itemize}

\begin{tcolorbox}[breakable,colback=teal!4,colframe=teal!35!black,title={Before you move on}]
What does \sk{दाडिमम्} mean? (\textquotedblleft{}A pomegranate\textquotedblright{}.) Say it once more.
\end{tcolorbox}
\section[\texorpdfstring{\sk{द्राक्षा} (drākṣā)}{द्राक्षा (drākṣā)}]{\sk{द्राक्षा} (drākṣā) — grapes}

\label{lesson:SA-C109-draksa}



\begin{quote}
Before the new one: say the Sanskrit for a coconut, then the Sanskrit for a pomegranate.
\end{quote}

\subsection*{You'll want to know: \sk{द्राक्षा}}
\textbf{\sk{द्राक्षा}} — \emph{drākṣā} — \textquotedblleft{}grapes\textquotedblright{}.

\begin{grammarlens}[title={What you've built}]
One more word for this chapter.
\end{grammarlens}

\subsection*{Your turn}
\begin{itemize}
\raggedright
  \item \emph{Say it:} \emph{drākṣā}
  \item \emph{Say it:} \emph{drākṣā}, once more
  \item \emph{Say it:} say \emph{drākṣā}
  \item \emph{Recall:} say the Sanskrit for a coconut, then the Sanskrit for a pomegranate, then say \emph{drākṣā} again
\end{itemize}

\begin{tcolorbox}[breakable,colback=teal!4,colframe=teal!35!black,title={Before you move on}]
What does \sk{द्राक्षा} mean? (\textquotedblleft{}Grapes\textquotedblright{}.) Say it once more.
\end{tcolorbox}
\section[\texorpdfstring{\sk{पनसम्} (panasam)}{पनसम् (panasam)}]{\sk{पनसम्} (panasam) — a jackfruit}

\label{lesson:SA-C109-panasam}



\begin{quote}
Before the new one: say the Sanskrit for a pomegranate, then the Sanskrit for grapes.
\end{quote}

\subsection*{You'll want to know: \sk{पनसम्}}
\textbf{\sk{पनसम्}} — \emph{panasam} — \textquotedblleft{}a jackfruit\textquotedblright{}.

\begin{grammarlens}[title={What you've built}]
One more word for this chapter.
\end{grammarlens}

\subsection*{Your turn}
\begin{itemize}
\raggedright
  \item \emph{Say it:} \emph{panasam}
  \item \emph{Say it:} \emph{panasam}, once more
  \item \emph{Say it:} say \emph{panasam}
  \item \emph{Recall:} say the Sanskrit for a pomegranate, then the Sanskrit for grapes, then say \emph{panasam} again
\end{itemize}

\begin{tcolorbox}[breakable,colback=teal!4,colframe=teal!35!black,title={Before you move on}]
What does \sk{पनसम्} mean? (\textquotedblleft{}A jackfruit\textquotedblright{}.) Say it once more.
\end{tcolorbox}
\section[\texorpdfstring{\sk{निम्बुकम्} (nimbukam)}{निम्बुकम् (nimbukam)}]{\sk{निम्बुकम्} (nimbukam) — a lemon}

\label{lesson:SA-C109-nimbukam}



\begin{quote}
Before the new one: say the Sanskrit for grapes, then the Sanskrit for a jackfruit.
\end{quote}

\subsection*{You'll want to know: \sk{निम्बुकम्}}
\textbf{\sk{निम्बुकम्}} — \emph{nimbukam} — \textquotedblleft{}a lemon\textquotedblright{}.

\begin{grammarlens}[title={What you've built}]
One more word for this chapter.
\end{grammarlens}

\subsection*{Your turn}
\begin{itemize}
\raggedright
  \item \emph{Say it:} \emph{nimbukam}
  \item \emph{Say it:} \emph{nimbukam}, once more
  \item \emph{Say it:} say \emph{nimbukam}
  \item \emph{Recall:} say the Sanskrit for grapes, then the Sanskrit for a jackfruit, then say \emph{nimbukam} again
\end{itemize}

\begin{tcolorbox}[breakable,colback=teal!4,colframe=teal!35!black,title={Before you move on}]
What does \sk{निम्बुकम्} mean? (\textquotedblleft{}A lemon\textquotedblright{}.) Say it once more.
\end{tcolorbox}
\section[\texorpdfstring{\sk{अपूपः} (apūpaḥ)}{अपूपः (apūpaḥ)}]{\sk{अपूपः} (apūpaḥ) — a cake}

\label{lesson:SA-C109-apupah}



\begin{quote}
Before the new one: say the Sanskrit for a jackfruit, then the Sanskrit for a lemon.
\end{quote}

\subsection*{You'll want to know: \sk{अपूपः}}
\textbf{\sk{अपूपः}} — \emph{apūpaḥ} — \textquotedblleft{}a cake\textquotedblright{}.

\begin{grammarlens}[title={What you've built}]
One more word for this chapter.
\end{grammarlens}

\subsection*{Your turn}
\begin{itemize}
\raggedright
  \item \emph{Say it:} \emph{apūpaḥ}
  \item \emph{Say it:} \emph{apūpaḥ}, once more
  \item \emph{Say it:} say \emph{apūpaḥ}
  \item \emph{Recall:} say the Sanskrit for a jackfruit, then the Sanskrit for a lemon, then say \emph{apūpaḥ} again
\end{itemize}

\begin{tcolorbox}[breakable,colback=teal!4,colframe=teal!35!black,title={Before you move on}]
What does \sk{अपूपः} mean? (\textquotedblleft{}A cake\textquotedblright{}.) Say it once more.
\end{tcolorbox}
\section[(dialogue)]{Tools and things you use — a first pass}

\label{lesson:SA-R109-first-pass-tools-and-things-you-use}



\begin{quote}
Say the words for a lemon and a cake.
\end{quote}

\subsection*{Your turn}
Say these aloud:

\begin{itemize}
\raggedright
  \item a nail, a hammer, scissors, a needle, thread
  \item a comb, soap, a towel, glasses, a bag
  \item a plate, a cup, a spoon, a stove, a candle
  \item a ball, a toy, a ring, a bangle, a shawl
  \item a sari, a shoe, a turban, a shirt, a telephone
\end{itemize}

\begin{tcolorbox}[breakable,colback=teal!4,colframe=teal!35!black,title={Before you move on}]
A nail? (\textbf{\sk{कीलः}}, \emph{kīlaḥ}.) A hammer? (\textbf{\sk{मुद्गरः}}, \emph{mudgaraḥ}.) Scissors? (\textbf{\sk{कर्तरी}}, \emph{kartarī}.) A needle? (\textbf{\sk{सूचिः}}, \emph{sūciḥ}.) Thread? (\textbf{\sk{सूत्रम्}}, \emph{sūtram}.) A comb? (\textbf{\sk{कंकतिका}}, \emph{kaṅkatikā}.) Soap? (\textbf{\sk{फेनकम्}}, \emph{phenakam}.) A towel? (\textbf{\sk{प्रोञ्छनम्}}, \emph{proñchanam}.) Glasses? (\textbf{\sk{उपनेत्रम्}}, \emph{upanetram}.) A bag? (\textbf{\sk{स्यूतः}}, \emph{syūtaḥ}.) A plate? (\textbf{\sk{स्थाली}}, \emph{sthālī}.) A cup? (\textbf{\sk{चषकः}}, \emph{caṣakaḥ}.) A spoon? (\textbf{\sk{चमसः}}, \emph{camasaḥ}.) A stove? (\textbf{\sk{चुल्ली}}, \emph{cullī}.) A candle? (\textbf{\sk{वर्तिका}}, \emph{vartikā}.) A ball? (\textbf{\sk{कन्दुकः}}, \emph{kandukaḥ}.) A toy? (\textbf{\sk{क्रीडनकम्}}, \emph{krīḍanakam}.) A ring? (\textbf{\sk{अंगुलीयकम्}}, \emph{aṅgulīyakam}.) A bangle? (\textbf{\sk{कंकणम्}}, \emph{kaṅkaṇam}.) A shawl? (\textbf{\sk{उत्तरीयम्}}, \emph{uttarīyam}.) A sari? (\textbf{\sk{शाटिका}}, \emph{śāṭikā}.) A shoe? (\textbf{\sk{पादत्राणम्}}, \emph{pādatrāṇam}.) A turban? (\textbf{\sk{उष्णीषम्}}, \emph{uṣṇīṣam}.) A shirt? (\textbf{\sk{कञ्चुकः}}, \emph{kañcukaḥ}.) A telephone? (\textbf{\sk{दूरभाषः}}, \emph{dūrabhāṣaḥ}.)
\end{tcolorbox}
\section[(dialogue)]{Machines, the house and food — a second pass}

\label{lesson:SA-R109-second-pass-machines-the-house-and-food}



\begin{quote}
Say the words for a television and a computer.
\end{quote}

\subsection*{Your turn}
Say these aloud:

\begin{itemize}
\raggedright
  \item a television, a computer, a machine, a fan, a box
  \item a rope, wheat, rice, mung beans, dal
  \item jaggery, pepper, turmeric, ginger, garlic
  \item an onion, a potato, an aubergine, a pumpkin, a coconut
  \item a pomegranate, grapes, a jackfruit, a lemon, a cake
\end{itemize}

\begin{tcolorbox}[breakable,colback=teal!4,colframe=teal!35!black,title={Before you move on}]
A television? (\textbf{\sk{दूरदर्शनम्}}, \emph{dūradarśanam}.) A computer? (\textbf{\sk{संगणकम्}}, \emph{saṅgaṇakam}.) A machine? (\textbf{\sk{यन्त्रम्}}, \emph{yantram}.) A fan? (\textbf{\sk{व्यजनम्}}, \emph{vyajanam}.) A box? (\textbf{\sk{पेटिका}}, \emph{peṭikā}.) A rope? (\textbf{\sk{रज्जुः}}, \emph{rajjuḥ}.) Wheat? (\textbf{\sk{गोधूमः}}, \emph{godhūmaḥ}.) Rice? (\textbf{\sk{तण्डुलः}}, \emph{taṇḍulaḥ}.) Mung beans? (\textbf{\sk{मुद्गः}}, \emph{mudgaḥ}.) Dal? (\textbf{\sk{सूपः}}, \emph{sūpaḥ}.) Jaggery? (\textbf{\sk{गुडः}}, \emph{guḍaḥ}.) Pepper? (\textbf{\sk{मरीचम्}}, \emph{marīcam}.) Turmeric? (\textbf{\sk{हरिद्रा}}, \emph{haridrā}.) Ginger? (\textbf{\sk{आर्द्रकम्}}, \emph{ārdrakam}.) Garlic? (\textbf{\sk{लशुनम्}}, \emph{laśunam}.) An onion? (\textbf{\sk{पलाण्डुः}}, \emph{palāṇḍuḥ}.) A potato? (\textbf{\sk{आलुकम्}}, \emph{ālukam}.) An aubergine? (\textbf{\sk{वृन्ताकम्}}, \emph{vṛntākam}.) A pumpkin? (\textbf{\sk{कूष्माण्डम्}}, \emph{kūṣmāṇḍam}.) A coconut? (\textbf{\sk{नारिकेलम्}}, \emph{nārikelam}.) A pomegranate? (\textbf{\sk{दाडिमम्}}, \emph{dāḍimam}.) Grapes? (\textbf{\sk{द्राक्षा}}, \emph{drākṣā}.) A jackfruit? (\textbf{\sk{पनसम्}}, \emph{panasam}.) A lemon? (\textbf{\sk{निम्बुकम्}}, \emph{nimbukam}.) A cake? (\textbf{\sk{अपूपः}}, \emph{apūpaḥ}.)
\end{tcolorbox}
